% SEGMENT OLP-0488-S01
% Part: applied-modal-logics
% Chapter: epistemic-logic
% Section: bisimulations

\documentclass[../../../include/open-logic-section]{subfiles}

\begin{document}

\olfileid{aml}{el}{bsd}

\olsection{இருதிசை ஒப்புறவுகள்}

நமது அறிவுநிலை மொழியின் வெளிப்படுத்தும் ஆற்றலைப் பற்றி எஞ்சியிருக்கும் ஒரு
கேள்வி, மாதிரிகளுக்கும் அவற்றில் மெய்யாகும் வாய்பாடுகளுக்கும் இடையிலான உறவைப்
பற்றியது. சில வாய்பாடுகள் ஒரு சட்டகத்தில் செல்லுபடியாகும்போது, அச்சட்டகங்கள்
குறிப்பிட்ட பண்புகளை நிறைவுசெய்வதையும் அவை உறுதிசெய்யும் என்று நமது சட்டகப்
பொருத்த முடிவுகளிலிருந்து கண்டோம். ஆனால் எடுத்துக்காட்டாக, தற்சுட்டு அம்பு
உள்ள ஓர் உலகையும், ஒவ்வொரு உலகமும் அடுத்ததற்கு இட்டுச்செல்லும் முடிவிலா
உலகத் தொடரையும் வேறுபடுத்த நமது வாய்ப்புநிலை மொழி அனுமதிக்கிறதா? அதாவது,
இவ்விரு உலகங்களில் ஒன்றில் மட்டும் மெய்யாகக்கூடிய ஏதேனும் வாய்பாடு~$A$
உள்ளதா?

இருதிசை ஒப்புறவு என்பது தொடர்புமுறை மாதிரிகளுக்கு இடையே நாம் வரையறுக்கக்கூடிய
ஓர் உறவு; அவை விளைவில் ஒரே கட்டமைப்பைக் கொண்டுள்ளன என்று அது கூறுகிறது.
நாம் காணவிருப்பதுபோல், நமது அறிவுநிலை மொழியில் பிடித்துக்கொள்ளக்கூடிய
மாதிரிகளுக்கிடையிலான ஒரு நிகர்த்தன்மையை அது பிடித்துக்கொள்ளும்.

\begin{defn}[இருதிசை ஒப்புறவு]
$M_1 = \tuple{W_1, R_1, V_1}$, $M_2 = \tuple{W_2, R_2, V_2}$ ஆகியவை இரு
தொடர்புமுறை மாதிரிகளாக இருக்கட்டும். மேலும், $\mathcal{R} \subseteq W_1
\times W_2$ ஓர் ஈருறுப்பு உறவாக இருக்கட்டும். ஒவ்வொரு
$\tuple{w_1, w_2} \in \mathcal{R}$ க்கும் பின்வருவன மெய்யாகும்போது,
$\mathcal{R}$ ஓர் \emph{இருதிசை ஒப்புறவு} என்கிறோம்:

\begin{enumerate}
  \item எல்லா !!{propositional variable}s~$p$ க்கும் $w_1 \in V_1(p)$
  எனில், எனில் மட்டுமே, $w_2 \in V_2(p)$.
  % SOURCE CORRECTION TA-AML-004: the chapter's agent-symbol set is G.
  \item எல்லா முகவர்கள் $a \in G$ மற்றும் உலகங்கள் $v_1 \in W_1$ க்கும்,
  $R_{1_a} w_1 v_1$ எனில், $R_{2_a} w_2 v_2$ ஆகவும்
  $\tuple{v_1, v_2} \in \mathcal{R}$ ஆகவும் இருக்கும் ஏதோவொரு
  $v_2 \in W_2$ உள்ளது.
  \item எல்லா முகவர்கள் $a \in G$ மற்றும் உலகங்கள் $v_2 \in W_2$ க்கும்,
   $R_{2_a} w_2 v_2$ எனில், $R_{1_a} w_1 v_1$ ஆகவும்
   $\tuple{v_1, v_2} \in \mathcal{R}$ ஆகவும் இருக்கும் ஏதோவொரு
   $v_1 \in W_1$ உள்ளது.
\end{enumerate}

$w_1$, $w_2$ ஆகிய உலகங்களை இணைக்கும் ஓர் இருதிசை ஒப்புறவு $M_1$, $M_2$
ஆகியவற்றுக்கு இடையில் இருக்கும்போது, $\tuple{M_1, w_1} \leftrightarroweq
\tuple{M_2, w_2}$ என்றும் எழுதலாம்; $\tuple{M_1, w_1}$,
$\tuple{M_2, w_2}$ ஆகியவை \emph{இருதிசை ஒப்புடையவை} என்றும் கூறலாம்.
\end{defn}

இருதிசை ஒப்புறவிலுள்ள வெவ்வேறு கூறுகள் வெவ்வேறு விடயங்களை உறுதிசெய்கின்றன.
எல்லா !!{propositional variable}s மீதும் உடன்பாட்டை உறுதிசெய்வதால், இருதிசை
ஒப்புடைய உலகங்கள் ஒரே வாய்ப்புநிலையற்ற வாய்பாடுகளை நிறைவுசெய்யும் என்பதை
கூறு~1 உறுதிசெய்கிறது. முறையே ``முன்னோக்கு'', ``பின்னோக்கு'' என்று சில
நேரங்களில் அழைக்கப்படும் மற்ற இரு கூறுகளும், அணுகுறவுகள் ஒரே கட்டமைப்பைக்
கொண்டிருப்பதை உறுதிசெய்கின்றன.

\begin{thm}
$\tuple{M_1, w_1} \leftrightarroweq \tuple{M_2, w_2}$ எனில், ஒவ்வொரு
!!{formula}~$!A$ க்கும், $\mSat{M_1}{!A}[w_1]$ எனில், எனில் மட்டுமே,
$\mSat{M_2}{!A}[w_2]$.
\end{thm}

\begin{figure}
  \begin{center}
    \begin{tikzpicture}[modal]
      \node[world] (w1) {$w_1$};
      \node[world] (w2) [above left=of w1]{$w_2$};
      \node[world] (w3) [above right=of w1] {$w_3$};
      \draw[<->] (w1) to node {$a$} (w2);
      \draw[->,loop below] (w1) to node {$a$} (w1);
      \draw[<->] (w1) to [swap] node {$a$} (w3);
      \draw[->,loop above] (w2) to node {$a$} (w2) ;
      \draw[->,loop above] (w3) to node {$a$} (w3) ;

      \node[world](v1)[right of=w1, xshift=1.5in]{$v_1$};
      \node[world](v2)[above of=v1]{$v_2$};
      \draw[<->] (v1) to node {$a$} (v2);
      \draw[->,loop below] (v1) to node {$a$} (v1);
      \draw[->,loop above] (v2) to node {$a$} (v2) ;

      \draw[-,dotted] (w1) to (v1) ;
      \draw[-,dotted,bend left=45] (w2) to (v2) ;
      \draw[-,dotted,bend left=30] (w3) to (v2) ;
    \end{tikzpicture}
  \end{center}
  \caption{இரு இருதிசை ஒப்புடைய மாதிரிகள்.}
  \ollabel{fig:bisimilar}
\end{figure}

\olref{fig:bisimilar} இல் காட்டப்பட்டுள்ள இரு மாதிரிகளும் ஒன்றுக்கொன்று
முற்றிலும் ஒரே மாதிரியானவை அல்ல என்றாலும், உலகங்கள் $w_1$,~$v_1$ ஆகியவற்றை
இணைக்கும் ஓர் இருதிசை ஒப்புறவு உள்ளது. இவ்விருதிசை ஒப்புறவு $w_2$, $w_3$
ஆகிய இரண்டையும்~$v_2$ உடன் இணைக்கும்; நமது வாய்ப்புநிலை மொழியில்
வெளிப்படுத்தக்கூடிய எதுவும் அவற்றை உண்மையில் வேறுபடுத்த முடியாது என்பதே இதன்
கருத்து. ஆனால் $w_2$,~$w_3$ ஆகியவை வெவ்வேறு !!{propositional variable}s ஐ
நிறைவுசெய்தால் நிலைமை வேறுபடும்.

\end{document}
